\chapter{Conclusions and Future Work}
\label{chap:conclusions}

\section{Conclusions}
\label{sec:conclusions}

This thesis built a transparent, reproducible top-tagging pipeline and used
it to compare four classifiers --- a boosted decision tree on engineered
features, a jet-image CNN, the ParticleNet graph network and the Particle
Transformer --- on the public 14\,TeV Top-Quark Tagging Reference Dataset,
with a focus on data efficiency and probability calibration rather than on a
new record. Returning to the objectives of
Chapter~\ref{chap:introduction}, the study reached the following
conclusions.

\begin{enumerate}
    \item \textbf{The established ordering is reproduced.} Under identical
    preprocessing and a modest budget, the Particle Transformer is best
    (AUC $=0.978$, $1/\varepsilon_B\approx136$ at $\varepsilon_S=0.5$),
    followed by ParticleNet, the CNN and the BDT, matching the community
    benchmark.
    \item \textbf{Data efficiency separates the models.} The BDT is the most
    data-efficient and is the strongest model at the smallest training size,
    but it is overtaken by the deep models as data grows. The Particle
    Transformer in particular rises from the weakest model at $5\times10^{3}$
    jets to the strongest at $5\times10^{4}$, a clear crossover with direct
    practical implications for low-statistics analyses.
    \item \textbf{Ranking and calibration are decoupled.} The highest-AUC
    model (the Particle Transformer) is comparatively well-calibrated,
    whereas the jet-image CNN --- only third on AUC --- is the
    worst-calibrated by an order of magnitude. A single-parameter temperature
    rescaling restores good calibration for the BDT, ParticleNet and the
    transformer, but only partially corrects the CNN. A high AUC therefore
    does not by itself license the use of a model's score as a probability.
    \item \textbf{World-class benchmarks are reproducible on free hardware.}
    The entire study was run on a free cloud GPU using open data and open
    tools, with automatic checkpointing and resumption, demonstrating that
    such work is accessible to students in resource-constrained settings.
\end{enumerate}

\section{Limitations}
\label{sec:limitations}

The study used subsamples of up to $5\times10^{4}$ jets rather than the full
$1.2\times10^{6}$, so absolute performances are slightly below the published
saturation values; the trends, not the records, are the contribution. The
reference dataset contains no pile-up and uses a single detector
parametrisation, so detector robustness and high-luminosity behaviour are not
addressed. Each architecture was trained with a fixed, modest configuration
rather than an exhaustive hyper-parameter search, and four models, while
representative, are not exhaustive.

\section{Future Work}
\label{sec:future}

Several extensions follow naturally from this work:
\begin{enumerate}
    \item \textbf{Scale up.} Repeat the data-efficiency sweep up to the full
    $1.2\times10^{6}$ jets on a larger GPU to locate precisely where the
    transformer overtakes ParticleNet by a statistically significant margin.
    \item \textbf{Richer recalibration.} Compare temperature scaling with
    non-parametric isotonic regression, particularly for the CNN whose
    miscalibration temperature scaling does not fully remove.
    \item \textbf{Robustness.} Study calibration and discrimination under
    simulated pile-up and detector smearing, which is where calibration
    matters most for real analyses.
    \item \textbf{More architectures.} Add a Deep~Sets baseline and a
    Lorentz-equivariant network~\citep{gong2022lorentznet} to the comparison.
    \item \textbf{Uncertainty quantification.} Use deep ensembles or
    Monte-Carlo dropout to attach predictive uncertainties to the tagger
    scores.
\end{enumerate}

In summary, this project shows that asking \emph{how much data} and
\emph{how trustworthy the probabilities are}, rather than only \emph{how
high the AUC}, yields a more complete and more useful picture of deep
jet-tagging models --- and that this picture can be obtained reproducibly on
freely available hardware.
